% paper/sections/03_method.tex
\section{Method}
\label{sec:method}

\subsection{Problem Formulation}

We model LLM routing as a contextual bandit problem. At each time step $t$:
\begin{enumerate}
    \item A query $x_t \in \mathcal{X}$ arrives (e.g., a natural language prompt).
    \item The router selects a model arm $a_t \in \mathcal{A} = \{1, \ldots, K\}$.
    \item The router receives a reward $r_t = R(a_t, x_t) \in [0, 1]$.
\end{enumerate}

The reward function balances quality and cost:
\begin{equation}
R_t = \alpha_q \cdot \text{acc}(a_t, x_t) - \alpha_c \cdot \frac{c_{a_t}}{c_{\max}}
\end{equation}
where $\alpha_q + \alpha_c = 1$, $c_{a_t}$ is the per-query cost of arm $a_t$, and $c_{\max}$ is the cost of the most expensive arm. For Experiment 1 (binary routing), $\alpha_q = 0.85$, $\alpha_c = 0.15$. For Experiment 2 (quality-dominant), $\alpha_q = 0.85$.

\subsection{Feature Space}
\label{sec:features}

We extract a 27-dimensional context vector $\phi(x_t) \in \mathbb{R}^{27}$ from each query:
\begin{itemize}
    \item \textbf{Message features [0--10]:} Log character count, log word count, sentence count, question mark count, number count, long-query indicator, code block pairs, has-question flag, capitalization flag, vocabulary richness, word saturation.
    \item \textbf{Extended features [11--16]:} Math/reasoning/code/creative keyword scores, difficulty estimate, vocabulary richness (alias).
    \item \textbf{Model features [17--26]:} Ten zero-initialized dimensions reserved for model-specific embeddings; populated from model metadata when available, otherwise zero-padded.
\end{itemize}
All features are L2-normalized before input: $\tilde{\phi}(x_t) = \phi(x_t) / \|\phi(x_t)\|$.

\subsection{LinUCB-27d}
\label{sec:linucb}

LinUCB \citep{li2010contextual} maintains per-arm ridge regression models. For arm $a$:
\begin{align}
A_a &\leftarrow A_a + \tilde{\phi}(x_t)\tilde{\phi}(x_t)^\top, \quad
b_a \leftarrow b_a + r_t \cdot \tilde{\phi}(x_t) \\
\hat{\theta}_a &= A_a^{-1} b_a, \quad
p_a(x_t) = \hat{\theta}_a^\top \tilde{\phi}(x_t) + \beta \sqrt{\tilde{\phi}(x_t)^\top A_a^{-1} \tilde{\phi}(x_t)}
\end{align}
The UCB score $p_a(x_t)$ trades off estimated quality (first term) with uncertainty (second term, controlled by $\beta$).

\subsection{LinTS-27d (Adapted from Abeille \& Lazaric, 2017)}
\label{sec:lints}

We adapt Linear Thompson Sampling \citep{abeille2017linear} for LLM routing. While Thompson Sampling for linear bandits is well-established\mbox{---}dating to \citet{agrawal2013thompson} and formalized by \citet{abeille2017linear}\mbox{---}its application to LLM routing with a 27-dimensional linguistic feature space is novel. For each arm $a$, LinTS maintains a Gaussian posterior over the weight vector $\theta_a$, given the observed data $\mathcal{D}_a = \{(\tilde{\phi}(x_i), r_i) : a_i = a, i \leq t\}$:
\begin{equation}
p(\theta_a \mid \mathcal{D}_a) = \mathcal{N}(\mu_a, v^2 \Sigma_a)
\end{equation}
where $\Sigma_a = A_a^{-1}$, $\mu_a = \Sigma_a b_a$. At each step, we sample:
\begin{equation}
\tilde{\theta}_a \sim \mathcal{N}(\mu_a, v^2 \Sigma_a) \quad \forall a \in \mathcal{A}
\end{equation}
and select $a^* = \arg\max_{a} \tilde{\phi}(x_t)^\top \tilde{\theta}_a$.

\textbf{Key advantage over LinUCB:} LinTS has no $\beta$ hyperparameter. Exploration is proportional to posterior variance, which decays naturally as data accumulates. We show in \S\ref{sec:ablations} that LinUCB's performance is sensitive to $\beta$, while LinTS requires no tuning.

\textbf{Regret bound:} LinTS achieves $O(d\sqrt{T \log T})$ cumulative regret with high probability \citep{abeille2017linear}, the same order as LinUCB.

\subsection{Relation to Full RouteSmith System}
\label{sec:product}

This paper focuses on RouteSmith's online contextual bandit router. The full RouteSmith product \citep{routesmith2025system} includes additional components not evaluated here:
\begin{itemize}
    \item \textbf{Random forest quality predictor:} A supervised router using RoRF-style features, following a cold-start $\rightarrow$ warm-up $\rightarrow$ learned phase progression. This serves as the default routing strategy in production when offline calibration data is available.
    \item \textbf{35-dimensional production features:} The production LinTS predictor (\texttt{LinTSPredictor}) defaults to 35 dimensions, extending the 27-dim benchmark features with additional model-metadata dimensions. The 27-dim features used in this paper represent a controlled subset for reproducible benchmarking.
    \item \textbf{Semantic cache:} Embedding-based response caching for semantically similar queries, orthogonal to the routing component.
    \item \textbf{Feedback collection:} Production telemetry for continuous routing improvement.
\end{itemize}
The LinTS bandit router described in this paper is one of multiple routing strategies available in the RouteSmith system, suitable for zero-shot deployment scenarios where pretraining data is unavailable.

\subsection{TS-Cat Baseline}
\label{sec:tscat}

As a context-free bandit baseline, TS-Cat maintains a Beta($\alpha_k$, $\beta_k$) posterior per query category $k$:
\begin{equation}
\theta_k \sim \text{Beta}(\alpha_k, \beta_k), \quad \alpha_k, \beta_k \leftarrow 1.0 \text{ (uniform prior)}
\end{equation}
Route to strong if $\theta_k > 0.5$. Update: if weak model answers correctly, $\beta_k \mathrel{+}= 1$; else $\alpha_k \mathrel{+}= 1$. TS-Cat captures category-level routing preferences but ignores query-level features.
